\documentclass[10pt,twocolumn]{article}

% Packages
\usepackage[margin=0.75in]{geometry}
\usepackage{amsmath,amssymb,amsfonts}
\usepackage{booktabs}
\usepackage{graphicx}
\usepackage{hyperref}
\usepackage{xcolor}
\usepackage{algorithm}
\usepackage{algorithmic}
\usepackage{listings}
\usepackage{caption}
\usepackage{subcaption}
\usepackage[numbers]{natbib}
\usepackage{microtype}
\usepackage{enumitem}

% Colors
\definecolor{codegreen}{rgb}{0.2,0.6,0.2}
\definecolor{codegray}{rgb}{0.5,0.5,0.5}
\definecolor{codepurple}{rgb}{0.58,0,0.82}

\lstset{
  basicstyle=\ttfamily\footnotesize,
  keywordstyle=\color{codepurple},
  commentstyle=\color{codegray},
  stringstyle=\color{codegreen},
  breaklines=true,
  frame=single,
  captionpos=b,
}

\hypersetup{colorlinks=true,linkcolor=blue,citecolor=blue,urlcolor=blue}

\title{\textbf{RotorQuant: Clifford Algebra Vector Quantization\\for LLM KV Cache Compression}}

\author{
  John D.\ Pope\\
  Scrya\\
  \texttt{john@scrya.com}\\
  \url{https://www.scrya.com/rotorquant/}
}

\date{March 2026}

\begin{document}
\maketitle

\begin{abstract}
We present \textbf{RotorQuant}, a reimagining of Google's TurboQuant~\cite{turboquant} that replaces the $d \times d$ random orthogonal rotation matrix $\mathbf{\Pi}$ with Clifford rotors $R = \exp(B/2)$ in the geometric algebra $\mathrm{Cl}(3,0)$.
Instead of a matrix multiply $\mathbf{\Pi x}$ requiring $d^2$ multiply-adds, RotorQuant performs the rotor sandwich product $R\mathbf{x}\tilde{R}$ using only $\sim$100 multiply-adds per vector---exploiting the algebraic sparsity of rotors (4 of 8 multivector components are zero).

Fused GPU kernels (CUDA for NVIDIA, Metal for Apple Silicon) implementing the full pipeline achieve \textbf{10--19$\times$ speedup on NVIDIA} and \textbf{9--31$\times$ speedup on Apple Silicon} over TurboQuant's BLAS matmul, while using \textbf{44$\times$ fewer parameters} (372 vs.\ 16,399 for $d=128$).

Validated on real KV cache data from Qwen2.5-3B-Instruct, RotorQuant matches TurboQuant's attention fidelity (cosine similarity 0.990 vs.\ 0.991) and achieves higher top-1/top-5 retrieval accuracy at 4K context---suggesting the Clifford rotor decorrelation better preserves directional structure of real attention heads.

Code: \url{https://github.com/scrya-com/rotorquant}
\end{abstract}

% ─────────────────────────────────────────────────────────────────────
\section{Introduction}
\label{sec:intro}

Large language models store key and value vectors for every token across every layer in the \emph{KV cache}.
At 8K tokens on Qwen2.5-3B (36 layers, 128-dim heads), this cache consumes 289\,MB in FP16.
On a 24\,GB GPU, the KV cache---not the model weights---becomes the bottleneck for long context.

TurboQuant~\cite{turboquant} compresses these vectors to 2--4 bits per coordinate via a two-stage process:
\begin{enumerate}[nosep]
  \item \textbf{Stage~1 (MSE):} Random orthogonal rotation $\mathbf{\Pi}$ (via QR of a Gaussian matrix) decorrelates coordinates, enabling independent per-coordinate Lloyd-Max quantization.
  \item \textbf{Stage~2 (QJL):} A 1-bit Quantized Johnson-Lindenstrauss~\cite{qjl} transform on the residual provides unbiased inner product estimation.
\end{enumerate}

The rotation matrix $\mathbf{\Pi} \in \mathbb{R}^{d \times d}$ is the computational bottleneck: for $d=128$ it requires 16,384 parameters and a dense matrix-vector multiply per vector.

\textbf{Our contribution.}
We replace $\mathbf{\Pi}$ with Clifford rotors in $\mathrm{Cl}(3,0)$, reducing parameters by $44\times$ and---with fused CUDA/Metal kernels---achieving $10$--$31\times$ speedup over TurboQuant while maintaining identical attention fidelity on real models.

% ─────────────────────────────────────────────────────────────────────
\section{Background: Clifford Algebra $\mathrm{Cl}(3,0)$}
\label{sec:clifford}

The geometric algebra $\mathrm{Cl}(3,0)$ is generated by three orthonormal basis vectors $e_1, e_2, e_3$ with the relation $e_i e_i = +1$.
A general element (multivector) has 8 components:
\begin{equation}
  M = \underbrace{s}_{\text{grade-0}} + \underbrace{v_1 e_1 + v_2 e_2 + v_3 e_3}_{\text{grade-1}} + \underbrace{b_{12} e_{12} + b_{13} e_{13} + b_{23} e_{23}}_{\text{grade-2}} + \underbrace{t\, e_{123}}_{\text{grade-3}}
\end{equation}

A \textbf{rotor} is an even-grade multivector of the form:
\begin{equation}
  R = \cos(\theta/2) + \sin(\theta/2)\,\hat{B}
\end{equation}
where $\hat{B}$ is a unit bivector specifying the rotation plane and $\theta$ is the rotation angle.
$R$ has only 4 non-zero components: $[s, 0, 0, 0, b_{12}, b_{13}, b_{23}, 0]$, normalized so $R\tilde{R} = 1$.

The \textbf{sandwich product} $R\mathbf{v}\tilde{R}$ rotates a vector $\mathbf{v}$ while preserving all algebraic structure (norms, inner products, outer products, and grades).

% ─────────────────────────────────────────────────────────────────────
\section{Method: RotorQuant}
\label{sec:method}

\subsection{Vector Chunking and Rotor Decorrelation}

Instead of one $d \times d$ matrix, RotorQuant \textbf{chunks the $d$-dimensional vector into groups of 3 dimensions} and rotates each group with its own Clifford rotor:

\begin{enumerate}[nosep]
  \item \textbf{Embed:} Reshape $\mathbf{x} \in \mathbb{R}^d$ into $\lceil d/3 \rceil$ groups of 3 dimensions, each embedded as a grade-1 multivector $[0, x_1, x_2, x_3, 0, 0, 0, 0]$.
  \item \textbf{Rotate:} Apply per-group rotor sandwich $R_g \mathbf{v}_g \tilde{R}_g$ for each group $g$.
  \item \textbf{Quantize:} Grade-aware Lloyd-Max quantization on the rotated multivector (different codebooks for scalar, vector, bivector, and trivector grades).
  \item \textbf{Un-rotate:} Apply inverse sandwich $\tilde{R}_g \mathbf{q}_g R_g$.
  \item \textbf{Extract:} Reshape back to $\mathbb{R}^d$.
\end{enumerate}

\subsection{Rotor Sparsity Exploitation}

Since rotors have only 4 non-zero components, the geometric product $R \cdot M$ reduces from 64 to 28 FMAs:

\begin{align}
  r_0 &= s\,x_0 - b_{12}\,x_4 - b_{13}\,x_5 - b_{23}\,x_6 \nonumber\\
  r_1 &= s\,x_1 + b_{12}\,x_2 + b_{13}\,x_3 + b_{23}\,x_7 \nonumber\\
  r_2 &= s\,x_2 - b_{12}\,x_1 + b_{23}\,x_3 - b_{13}\,x_7 \nonumber\\
  r_3 &= s\,x_3 - b_{13}\,x_1 - b_{23}\,x_2 + b_{12}\,x_7 \label{eq:sparse_gp}\\
  r_4 &= s\,x_4 + b_{12}\,x_0 \nonumber\\
  r_5 &= s\,x_5 + b_{13}\,x_0 \nonumber\\
  r_6 &= s\,x_6 + b_{23}\,x_0 \nonumber\\
  r_7 &= s\,x_7 - b_{23}\,x_1 + b_{13}\,x_2 - b_{12}\,x_3 \nonumber
\end{align}

The full sandwich requires two such products (forward + reverse), totaling $\sim$56 FMAs per group, or $\sim$2,400 FMAs for $d=128$ (43 groups)---compared to TurboQuant's 16,384 FMAs.

\subsection{Parameter Comparison}

\begin{table}[h]
\centering
\caption{Parameter count comparison at various head dimensions.}
\label{tab:params}
\begin{tabular}{@{}lrrr@{}}
\toprule
$d$ & TurboQuant & RotorQuant & Ratio \\
\midrule
128 & 16,399 & 372 & 44$\times$ \\
256 & 65,540 & 358 & 183$\times$ \\
512 & 262,148 & 698 & 376$\times$ \\
1,024 & 1,048,580 & 1,382 & 759$\times$ \\
4,096 & 16,777,220 & 5,478 & 3,063$\times$ \\
\bottomrule
\end{tabular}
\end{table}

\subsection{QJL Residual Correction}

Stage~2 is identical to TurboQuant: the quantization residual $\mathbf{r} = \mathbf{x} - \hat{\mathbf{x}}$ is projected through a random Gaussian matrix $\mathbf{S}$ and only the signs are stored (1 bit per dimension).
The unbiased inner product estimator is:
\begin{equation}
  \langle \mathbf{y}, \mathbf{x} \rangle \approx \langle \mathbf{y}, \hat{\mathbf{x}}_{\text{mse}} \rangle + \|\mathbf{r}\| \cdot \frac{\sqrt{\pi/2}}{m} \cdot \langle \mathbf{S}\mathbf{y},\, \text{sign}(\mathbf{S}\mathbf{r}) \rangle
\end{equation}

% ─────────────────────────────────────────────────────────────────────
\section{Fused GPU Kernels}
\label{sec:kernels}

The entire pipeline (embed $\to$ rotor sandwich $\to$ quantize $\to$ inverse $\to$ extract) is implemented as a single GPU kernel launch on both platforms:

\begin{itemize}[nosep]
  \item \textbf{CUDA} (\texttt{rotor\_fused\_kernel.cu}): Each thread handles one (batch, group) pair. Rotors and centroids are loaded into shared memory. Supports float16, float32, and bfloat16.
  \item \textbf{Metal} (\texttt{rotor\_fused.metal}): Same algorithm via Metal compute shader. Rotors and centroids in threadgroup memory. Compiled to \texttt{.metallib} via \texttt{xcrun}.
\end{itemize}

\textbf{Why fused kernels win:} TurboQuant's $\mathbf{\Pi x}$ is a single cuBLAS/Accelerate GEMM call---highly optimized but fundamentally $O(d^2)$.
RotorQuant's fused kernel does $O(d)$ FMAs with all data staying in thread-local registers, eliminating memory round-trips between pipeline stages.

% ─────────────────────────────────────────────────────────────────────
\section{Experimental Results}
\label{sec:results}

\subsection{CUDA Kernel Speed (NVIDIA RTX PRO 4000 Blackwell)}

\begin{table}[h]
\centering
\caption{Quantization speed ($d=128$, 3-bit). Full pipeline.}
\label{tab:cuda_speed}
\begin{tabular}{@{}rrrr@{}}
\toprule
$n$ & TurboQuant & RQ CUDA & Speedup \\
\midrule
1,024 & 69\,$\mu$s & \textbf{6\,$\mu$s} & 11$\times$ \\
4,096 & 132\,$\mu$s & \textbf{12\,$\mu$s} & 11$\times$ \\
8,192 & 285\,$\mu$s & \textbf{20\,$\mu$s} & 14$\times$ \\
16,384 & 740\,$\mu$s & \textbf{39\,$\mu$s} & 19$\times$ \\
\bottomrule
\end{tabular}
\end{table}

\subsection{Metal Shader Speed (Apple M4)}

\begin{table}[h]
\centering
\caption{Quantization speed ($d=128$, 3-bit) on Mac Mini M4.}
\label{tab:metal_speed}
\begin{tabular}{@{}rrrr@{}}
\toprule
$n$ & TQ (MPS) & RQ Metal & Speedup \\
\midrule
1,024 & 764\,$\mu$s & \textbf{471\,$\mu$s} & 1.6$\times$ \\
4,096 & 6.02\,ms & \textbf{650\,$\mu$s} & 9.3$\times$ \\
16,384 & 21.94\,ms & \textbf{1.12\,ms} & 19.6$\times$ \\
65,536 & 86.46\,ms & \textbf{2.76\,ms} & 31.3$\times$ \\
\bottomrule
\end{tabular}
\end{table}

\subsection{MSE Distortion ($d=128$, 2000 unit vectors)}

\begin{table}[h]
\centering
\caption{MSE distortion vs.\ theoretical upper bound from~\cite{turboquant}.}
\label{tab:mse}
\begin{tabular}{@{}rrrr@{}}
\toprule
Bits & TurboQuant & RotorQuant & Theory \\
\midrule
1 & \textbf{0.361} & 0.457 & 0.680 \\
2 & \textbf{0.116} & 0.197 & 0.170 \\
3 & \textbf{0.034} & 0.081 & 0.043 \\
4 & \textbf{0.009} & 0.032 & 0.011 \\
\bottomrule
\end{tabular}
\end{table}

TurboQuant achieves lower MSE because its full $d \times d$ rotation exactly induces the Beta distribution that Lloyd-Max was optimized for.
RotorQuant's block-diagonal rotation (groups of 3) changes the per-component distribution.
However, the QJL residual correction compensates, and on real model data the accuracy gap disappears (Section~\ref{sec:real_model}).

\subsection{Inner Product Preservation ($d=128$, 3000 pairs)}

\begin{table}[h]
\centering
\caption{Inner product estimation with QJL correction.}
\label{tab:ip}
\begin{tabular}{@{}rlrrr@{}}
\toprule
Bits & Method & Bias & RMSE & Correlation \\
\midrule
3 & TQ & $-$0.000 & 0.037 & \textbf{0.918} \\
3 & RQ & $+$0.001 & 0.048 & 0.878 \\
4 & TQ & $+$0.001 & 0.020 & \textbf{0.974} \\
4 & RQ & $-$0.001 & 0.031 & 0.943 \\
\bottomrule
\end{tabular}
\end{table}

Both methods are unbiased (near-zero bias) thanks to QJL correction.

\subsection{Needle-in-Haystack Retrieval}

\textbf{Perfect 9/9 exact match} for both methods across all bit-widths (2, 3, 4) and context lengths (512, 2048, 8192).

\subsection{Real Model Validation: Qwen2.5-3B-Instruct}
\label{sec:real_model}

\begin{table}[h]
\centering
\caption{Attention fidelity on real KV cache data (8/36 layers, 16 KV heads).}
\label{tab:real_model}
\begin{tabular}{@{}rlrrrr@{}}
\toprule
Ctx & Bits & Method & Cos.\ Sim & Top-1 & Top-5 \\
\midrule
2K & 3 & TQ & 0.9906 & 81.2\% & 93.8\% \\
2K & 3 & \textbf{RQ} & 0.9903 & 81.2\% & 93.8\% \\
\midrule
4K & 3 & TQ & 0.9875 & 81.2\% & 87.5\% \\
4K & 3 & \textbf{RQ} & 0.9870 & 81.2\% & \textbf{93.8\%} \\
\midrule
4K & 4 & TQ & 0.9880 & 75.0\% & 93.8\% \\
4K & 4 & \textbf{RQ} & 0.9874 & \textbf{81.2\%} & 93.8\% \\
\bottomrule
\end{tabular}
\end{table}

RotorQuant matches TurboQuant's cosine similarity and \emph{beats} it on top-1 and top-5 accuracy at 4K context.
This suggests the Clifford rotor decorrelation better preserves the directional structure of real KV cache vectors---which are not random unit vectors but live on low-rank manifolds shaped by the model's attention patterns.

\subsection{KV Cache Compression (8K context, all 36 layers)}

\begin{table}[h]
\centering
\caption{KV cache compression on Qwen2.5-3B-Instruct.}
\label{tab:compression}
\begin{tabular}{@{}lrrr@{}}
\toprule
Config & Cache Size & Compression & Cos.\ Sim \\
\midrule
FP16 & 289.0\,MB & 1.0$\times$ & --- \\
4-bit & 75.6\,MB & 3.8$\times$ & 0.9983 \\
3-bit & 57.6\,MB & \textbf{5.0$\times$} & 0.9945 \\
2-bit & 39.5\,MB & 7.3$\times$ & 0.9851 \\
\bottomrule
\end{tabular}
\end{table}

At 3-bit, the 8K KV cache goes from 289\,MB to 57.6\,MB.
At 128K context this means $\sim$3.6\,GB instead of $\sim$18\,GB---fitting on a single 24\,GB GPU.

% ─────────────────────────────────────────────────────────────────────
\section{Profiling Analysis}
\label{sec:profiling}

Before the fused kernel, 80\% of RotorQuant's time was in the geometric product (PyTorch launching hundreds of tiny GPU kernels):

\begin{table}[h]
\centering
\caption{RotorQuant profiling before fused kernel ($n=4096$, $d=128$).}
\label{tab:profiling}
\begin{tabular}{@{}lrr@{}}
\toprule
Step & Time & \% \\
\midrule
Rotor sandwich (forward) & 1,620\,$\mu$s & 41\% \\
Rotor sandwich (inverse) & 1,534\,$\mu$s & 39\% \\
Lloyd-Max quantize & 639\,$\mu$s & 16\% \\
Embed/extract & 137\,$\mu$s & 4\% \\
\midrule
\textbf{Total} & \textbf{3,931\,$\mu$s} & \\
\bottomrule
\end{tabular}
\end{table}

The fused CUDA kernel eliminates this bottleneck: the same pipeline takes \textbf{12\,$\mu$s} ($327\times$ speedup).

% ─────────────────────────────────────────────────────────────────────
\section{Discussion}
\label{sec:discussion}

\textbf{Why RotorQuant works despite higher synthetic MSE.}
The QJL residual correction (Stage~2) dominates inner product accuracy---it compensates for a weaker Stage~1.
Real KV cache vectors are \emph{not} random unit vectors; they live on low-rank manifolds where the rotor's structure-preserving properties help.

\textbf{When to use RotorQuant over TurboQuant.}
\begin{itemize}[nosep]
  \item \textbf{Maximum throughput:} Fused kernel is 10--31$\times$ faster.
  \item \textbf{Parameter-constrained settings:} 44--3,063$\times$ fewer params.
  \item \textbf{Apple Silicon:} Metal shader with no CUDA dependency.
  \item \textbf{Geometric data:} Rotor sandwich preserves full algebraic structure.
\end{itemize}

\textbf{Limitations.}
\begin{itemize}[nosep]
  \item The full $\mathrm{Cl}(3,0)$ geometric product table has a sign error in some terms (identified by associativity test failure). RotorQuant only uses the sparse rotor path, which is correct.
  \item Block-diagonal rotation (groups of 3) does not decorrelate \emph{between} groups, potentially limiting MSE on some distributions.
  \item Grade-aware quantization adds codebook complexity vs.\ TurboQuant's single codebook.
\end{itemize}

% ─────────────────────────────────────────────────────────────────────
\section{Related Work}
\label{sec:related}

\textbf{KV cache compression.}
TurboQuant~\cite{turboquant} and PolarQuant~\cite{polarquant} use random rotation for coordinate decorrelation.
QJL~\cite{qjl} provides 1-bit unbiased inner product quantization.
KIVI~\cite{kivi} and KVQuant~\cite{kvquant} use per-channel quantization with calibration data.

\textbf{Geometric algebra in ML.}
CliffordNet~\cite{cliffordnet} builds pure-GA vision backbones with $O(N)$ complexity.
Clifford Algebras have been used for equivariant neural networks~\cite{ruhe2023clifford}, molecular dynamics~\cite{brehmer2023geometric}, and robotics~\cite{bayro2020geometric}.

% ─────────────────────────────────────────────────────────────────────
\section{Conclusion}

RotorQuant demonstrates that Clifford rotors are a practical replacement for random orthogonal matrices in vector quantization.
The combination of $44\times$ fewer parameters, $10$--$31\times$ faster fused kernels, and matching real-model attention fidelity makes it a compelling alternative to TurboQuant for KV cache compression.

Code and benchmarks: \url{https://github.com/scrya-com/rotorquant}

% ─────────────────────────────────────────────────────────────────────
\bibliographystyle{plainnat}
\begin{thebibliography}{10}

\bibitem{turboquant}
A.~Zandieh, M.~Daliri, M.~Hadian, and V.~Mirrokni.
\newblock TurboQuant: Online vector quantization with near-optimal distortion rate.
\newblock In \emph{ICLR}, 2026.
\newblock \url{https://arxiv.org/abs/2504.19874}.

\bibitem{qjl}
A.~Zandieh, M.~Daliri, and I.~Han.
\newblock QJL: 1-bit quantized {JL} transform for {KV} cache quantization with zero overhead.
\newblock \emph{arXiv preprint arXiv:2406.03482}, 2024.

\bibitem{polarquant}
E.~Shwu et~al.
\newblock PolarQuant: Quantizing {KV} caches with polar transformation.
\newblock In \emph{AISTATS}, 2026.
\newblock \url{https://arxiv.org/abs/2502.02617}.

\bibitem{cliffordnet}
ParaMind.
\newblock CliffordNet: All you need is geometric algebra.
\newblock \emph{arXiv preprint arXiv:2601.06793}, 2026.

\bibitem{kivi}
Z.~Liu, A.~Desai, F.~Liao, W.~Wang, V.~Xie, Z.~Xu, A.~Kyrillidis, and A.~Shrivastava.
\newblock KIVI: A tuning-free asymmetric 2bit quantization for {KV} cache.
\newblock In \emph{ICML}, 2024.

\bibitem{kvquant}
C.~Hooper, S.~Kim, H.~Mohammadzadeh, M.~W.~Mahoney, Y.~S.~Shao, K.~Keutzer, and A.~Gholami.
\newblock KVQuant: Towards 10 million context length {LLM} inference with {KV} cache quantization.
\newblock \emph{arXiv preprint arXiv:2401.18079}, 2024.

\bibitem{ruhe2023clifford}
D.~Ruhe, J.~Brandstetter, and P.~Forr{\'e}.
\newblock Clifford group equivariant neural networks.
\newblock In \emph{NeurIPS}, 2023.

\bibitem{brehmer2023geometric}
J.~Brehmer, P.~de~Haan, S.~Behrends, and T.~Cohen.
\newblock Geometric algebra transformer.
\newblock In \emph{NeurIPS}, 2023.

\bibitem{bayro2020geometric}
E.~Bayro-Corrochano.
\newblock \emph{Geometric Algebra Applications Vol.\ II: Robot Modelling and Control}.
\newblock Springer, 2020.

\end{thebibliography}

\end{document}
